% M16, 11.09.2026: Durchstich am Einstieg erläutert; Ausgangsziel erhalten.
% M07, 11.09.2026: fachlicher Szenariorahmen ergänzt; methodische
% Wiederholungen verdichtet. Zielherkunft und E-ID-Funktion erhalten.
% Quelle des Szenariorahmens: input/transcripts/Interview-Einrichtung.txt,
% Gesprächsauftakt und Vorschlag einer App zur Kommunikationsunterstützung.

\chapter{Explorative Problemklärung und Herleitung der
Designanforderungen}
\label{chap:explorative-problemklaerung}

Auf Grundlage des Vorgehens aus
Abschnitt~\ref{sec:iteratives-vorgehen-evidenz} rekonstruiert dieses
Kapitel die architekturprägenden Entwicklungsbefunde. Aus ausgewählten
Beobachtungen, dem Nutzungsziel und technischen Konsequenzen werden
die zehn Designanforderungen A1 bis A10 hergeleitet.
Abschnitt~\ref{sec:konsolidierte-designanforderungen} führt sie
zusammen; Kapitel~\ref{chap:loesungsarchitektur} erläutert die daraus
entwickelte Architektur.

\section{Ausgangslage und Auswahl der Entwicklungsbefunde}
\label{sec:ausgangslage-entwicklungsbefunde}

Das als Projektmaterial verwendete Ausgangstranskript behandelt ein
Kick-off-Gespräch zur Planung einer App für eine Einrichtung für
Menschen mit Beeinträchtigungen.\footnote{Projektquelle im Repository:
\path{input/transcripts/Interview-Einrichtung.txt}.}
Darin erörtern die Einrichtungsleitung
und ein Projektteam aus Informatik und Sozialer Arbeit, wie Wissen
über betreute Personen und ihre individuelle Kommunikation zugänglich
gemacht werden könnte. Die vorgeschlagene App soll beispielsweise
Betreuenden helfen, Ausdrucksweisen einer Person anhand hinterlegter
Informationen besser zu verstehen. Diese App ist das zu planende
Zielsystem. Gegenstand der vorliegenden Arbeit ist das Agentensystem,
das solche Gesprächsinhalte für die Softwareplanung verarbeitet.

Der ursprüngliche technische Auftrag bestand im Weg vom Transkript
über KI-Agenten mit Microsoft Agent Framework zum GitHub-Issue
(Abschnitt~\ref{sec:iteratives-vorgehen-evidenz}). Das später um
Rückkopplung und Wiederverwendung erweiterte Nutzungsziel gab Anlass,
auch weitere Projektinformationen und den Umgang mit bereits
erzeugten Ergebnissen zu untersuchen. Die heutige Architektur war
damit noch nicht festgelegt.

Den Anfang bildete ein einfacher technischer Durchstich, also ein
erster lauffähiger Prototyp für den Weg von der Eingabe zum Ergebnis.
Ein einzelner Agent sollte das Transkript verarbeiten, Artefakte über
Dateiwerkzeuge ablegen und Freigaben anfordern. Die operative
Arbeitsumgebung wurde früh angebunden. Offen blieben der angemessene
Autonomiegrad, die Aufgabenverteilung, die Zustandsführung und die
menschliche Kontrolle. Die Entwicklung klärte schrittweise, welche
Entscheidungen dem Agenten überlassen werden konnten und wo explizite
Informationsführung, deterministische Prüfungen oder menschliche
Autorisierung benötigt wurden. Den Einsatz unterschiedlicher Modelle in frühen Entwicklungsständen
ordnet Abschnitt~\ref{sec:explorationsverlauf} methodisch ein.

Die Darstellung folgt den Übergängen, die eine tragende Anforderung,
Architekturentscheidung oder Grenze erklären; eine vollständige Chronik
ist nicht beabsichtigt. Abschnitt~\ref{sec:explorationsverlauf} gibt
zunächst den Überblick. Als formative Evidenz verwendete
Projektbeobachtungen erhalten lokale Evidenzanker (E-IDs).
Anhang~\ref{anh:evidenzanker} ordnet sie der Dokumentation im
eingefrorenen Repository-Stand und ihrer jeweiligen Aussagegrenze zu.
